% ============================================================================
%  CONCLUSIONES
% ============================================================================

\section{Conclusiones}

El proyecto de migración heterogénea de la base de datos de la Cooperativa de
Ahorro y Crédito Horizonte PyME S.R.L. desde MySQL 5.7 hacia PostgreSQL 16 se
ha ejecutado de forma íntegra, segura y verificable, cumpliendo con los ocho
objetivos específicos definidos en el Capítulo~II.

\begin{enumerate}
  \item \textbf{Diagnóstico del sistema de origen:} se levantó un inventario
        completo del motor MySQL 5.7.44: 8 tablas, 2 vistas, 19 entradas de
        índice, 0.28 MB de tamaño, un único usuario de aplicación
        (\texttt{app\_legacy\_user} con \texttt{GRANT ALL}), ausencia de claves
        foráneas y tráfico sin cifrado.
  \item \textbf{Incompatibilidades de tipos y semántica:} se documentó la matriz
        de conversión (\texttt{TINYINT(1)} $\rightarrow$ \texttt{BOOLEAN},
        \texttt{ENUM} $\rightarrow$ \texttt{VARCHAR} + \texttt{CHECK},
        \texttt{DATETIME} $\rightarrow$ \texttt{TIMESTAMPTZ} en
        \texttt{America/La\_Paz}, \texttt{INT UNSIGNED AUTO\_INCREMENT}
        $\rightarrow$ \texttt{BIGINT GENERATED ALWAYS AS IDENTITY}).
  \item \textbf{Diseño del esquema destino:} se implementaron tipos nativos de
        PostgreSQL 16, restricciones \texttt{CHECK}, 9 claves foráneas, 10
        índices y 3 vistas (incluyendo la vista enmascarada
        \texttt{vw\_socios\_audit}).
  \item \textbf{ETL determinista:} se reemplazó el ETL artesanal en Python por
        una arquitectura basada en \textbf{pgloader} con schema de staging
        (\texttt{stg}), precheck de validación, carga en una sola transacción con
        \texttt{OVERRIDING SYSTEM VALUE}, postcheck bidireccional
        (\texttt{to\_jsonb} + \texttt{EXCEPT ALL}), reajuste de secuencias con
        \texttt{setval} y \texttt{ANALYZE}. El proceso es idempotente y con
        rollback transaccional garantizado.
  \item \textbf{Endurecimiento de PostgreSQL 16:} TLS 1.2 obligatorio con
        \texttt{hostssl}, autenticación \texttt{scram-sha-256}, revocación de
        \texttt{PUBLIC}, 7 roles con mínimo privilegio
        (\texttt{admin\_dba}, \texttt{app\_core\_rw}, \texttt{auditor\_ro},
        \texttt{migration\_loader} y 4 roles de administradores del equipo),
        registro de conexiones, desconexiones y DDL.
  \item \textbf{Enmascaramiento de datos sensibles:} la vista
        \texttt{vw\_socios\_audit} expone CI/NIT, teléfono, correo y números de
        cuenta enmascarados; \texttt{auditor\_ro} no tiene acceso directo a
        tablas base sensibles.
  \item \textbf{Respaldo cifrado y restauración:} \texttt{pg\_dump -Fc} cifrado
        con AES-256-CBC y PBKDF2 (200.000 iteraciones), verificación SHA-256 y
        prueba de restauración en base temporal con los 55 socios recuperados.
  \item \textbf{Automatización y evidencias:} scripts reproducibles
        (\texttt{01\_audit\_origin.sh}, \texttt{02\_migrate\_and\_backup.sh},
        \texttt{03\_verify\_all.sh}) con salidas reales en \texttt{evidence/}.
\end{enumerate}

El proyecto cumple con los objetivos planteados: la base de datos de la
cooperativa ahora opera sobre PostgreSQL 16, un motor moderno y soportado, con
seguridad por capas, integridad referencial, datos sensibles enmascarados,
respaldo cifrado verificable y un proceso de migración reproducible y
automatizado.

\section{Recomendaciones}

A partir de la ejecución y de las pruebas realizadas se recomienda, para una
adopción en producción:

\begin{enumerate}
  \item \textbf{Gestión centralizada de secretos:} sustituir \texttt{.env} por un
        gestor (Vault, AWS Secrets Manager o equivalente), con rotación
        periódica de la frase de paso de los respaldos y separación entre la
        custodia de la clave y la del archivo cifrado.
  \item \textbf{Respaldo fuera del sitio:} aplicar la regla 3-2-1 y automatizar
        la verificación SHA-256 y la prueba de restauración periódica, no solo
        la generación del respaldo.
  \item \textbf{Reducción de privilegios en el origen:} sustituir el
        \texttt{GRANT ALL} de \texttt{app\_legacy\_user} por permisos mínimos de
        lectura antes de cualquier nueva corrida sobre datos reales.
  \item \textbf{Reproducibilidad de imágenes:} fijar las imágenes
        \texttt{mysql:5.7} y \texttt{postgres:16} por digest para evitar cambios
        silenciosos entre despliegues.
  \item \textbf{Monitoreo y capacidad:} instrumentar conexiones, consultas
        lentas y crecimiento de los datos; evaluar PgBouncer y, cuando el
        volumen lo justifique, replicación con CDC para reducir la ventana de
        corte.
  \item \textbf{Operación formal:} documentar un \emph{runbook} de reversión y
        ejecutar ensayos periódicos de recuperación con RTO y RPO medidos.
\end{enumerate}
